\subsection{Results and adoption decision}
All 90 workers pass their independent answer and full base-closure checks.
At SF0.1, Q5 candidate visits fall from 621,901 to 155,969 in every block.
Each lookahead invocation performs 30 extra probes, examines 16 physical rows,
retains two hints, and changes 101 ordering choices. The paired
lookahead/compiled join ratio is 0.509923, with observed range
$[0.480504,0.519889]$. Against \code{b1254c4}, which also differs in earlier
indexing and interpreter improvements, the join ratio is 0.273615
$[0.268051,0.284502]$. These comparisons isolate ordering only when the
denominator is the same-source compiled arm.

The result does not establish a robust pipeline gain. Large-Q5 component-total
paired ratios are 0.975196 $[0.971553,1.272566]$, and process ratios are
0.990329 $[0.976037,1.284827]$. Separate component-total medians are 9.298 seconds
with lookahead and 7.307 seconds compiled, despite the median paired ratio
being below one. The block-specific shifts explain why the two summary
statistics differ; neither should be substituted for the other. One block has
a slower lookahead pipeline, largely in phases preceding query execution.
All observations are retained. The experiment did not control desktop activity
or processor affinity, and the preparation code is shared between these arms.
The data do not establish that lookahead caused the preceding-phase variation.

The other five TPC-H components keep their candidate-row counts unchanged.
Their paired join ratios range from 0.996 to 1.043, with the larger Q9 case
about 4.3\% slower. The four adverse-control join ratios are 1.123 for unique
joins, 1.066 for nonselective joins, 1.155 for skew, and 1.238 for the empty
cycle. Thus the extra decision machinery has measurable costs even when it
does not change the traversal. The very large nonselective improvement against
\code{b1254c4} is already present in the compiled control and must not be
attributed to lookahead.

The heuristic remains disabled by default. Its conditional Q5 success supports
further investigation of intermediate-sensitive ordering, while the adverse
results reject a universal benefit. No workload-specific dispatcher was tuned
on these final observations. The stored protocol was frozen locally before
measurement, not independently timestamped or publicly preregistered. This
post-hoc discovery-workload result establishes neither RPT+/SYA superiority
nor worst-case-optimal or no-regret evaluation.
